\subsection{The remaining mixed thin systems by two residue obstructions}
\label{sec:DE-A3-mixed-thin-residues}

\begin{lemma}[The ordinary Frobenius algebra and its two pencils]
\label{lem:DE-A3-ordinary-two-pencils}
\status{Proved; root and reduction-descent full independent proof/source/code audits passed.}
Let $E$ be an ordinary elliptic curve over an algebraically closed
field of characteristic two, let $F:E\to E'=E^{(2)}$ be relative
Frobenius, and let $H=\OO_{E'}(qO)$, where $q\ge3$.
There is an orthogonal decomposition
\[
 F_*\OO_E=\OO_{E'}\oplus B,\qquad
 B\not\simeq\OO_{E'},\qquad
 \mu:B^{\otimes2}\xrightarrow{\sim}\OO_{E'}.
\]
Here orthogonality refers to the multiplication--Cartier pairing,
normalized to take the value one on $(1,1)$; in this decomposition
the pairing is $(a,b),(c,d)\mapsto ac+\mu bd$.
Put $V=F_*F^*H=H\oplus HB$.
One can choose $a_0\in H^0(H)$ and $b_2\in H^0(HB)$ with
reduced, disjoint zero divisors $A$ and $D$.
For every such choice the following maps are isomorphisms:
\[
\begin{split}
 H^0(H)\oplus H^0(HB)&\longrightarrow H^0(H^2),\\
 (x,y)&\longmapsto a_0x+\mu b_2y,\\
 H^0(HB)&\longrightarrow H^0(A,HB|_A),\\
 H^0(H)&\longrightarrow H^0(D,H|_D).
\end{split}
\]
\end{lemma}
\begin{proof}
Choose nonzero invariant differentials on $E,E'$.
Ordinarity says that Cartier on the invariant differential is nonzero.
After normalization it gives an $\OO_{E'}$-linear map
$\tau:F_*\OO_E\to\OO_{E'}$ with $\tau(1)=1$.
The pairing $(u,v)\mapsto\tau(uv)$ is perfect: this is also immediate
in a local $2$-basis from the Cartier formula.
For the original Cartier isomorphism and local formula see
\cite[0, Theorem~2.1.9, p.~515]{Illusie1979DRW}.
Thus $B=\ker\tau$ is a line bundle, orthogonal to the unit summand.
Every square in $F_*\OO_E$ belongs to the unit summand.
Locally a generator of $B$ therefore has square in $\OO_{E'}$,
and perfection of the pairing makes that square a unit.
This proves the stated multiplication isomorphism $\mu$.
Finally $H^0(F_*\OO_E)=k$ gives $H^0(B)=0$.
Hence $B$ is nontrivial of degree zero, and $H^1(B)=0$ as well.

The degree-$q$ line bundles $H$ and $HB$ are very ample, so general
sections have reduced disjoint zero divisors.
The first displayed map comes from the exact sequence
\[
 0\longrightarrow B
 \xrightarrow{\,t\mapsto(\mu b_2t,a_0t)\,}
 H\oplus HB
 \xrightarrow{\,(x,y)\mapsto a_0x+\mu b_2y\,}
 H^2\longrightarrow0.
\]
Its exactness follows from the absence of common zeros of $a_0,b_2$.
The vanishing of both cohomology groups of $B$ proves the first
isomorphism.  The two evaluation isomorphisms follow in the same way
from $HB(-A)=B$ and $H(-D)=B^{-1}=B$.
\end{proof}

\begin{lemma}[A nonzero stable-rank obstruction determinant]
\label{lem:DE-A3-mixed-residue-determinant}
\status{Proved; root and reduction-descent full independent proof/source/code audits passed.}
With the notation of Lemma~\ref{lem:DE-A3-ordinary-two-pencils},
there exist $a_1\in H^0(H)$ and $b_1\in H^0(HB)$ for which the
section triple
\[
 f_0=(a_0,0),\qquad f_1=(a_1,b_1),\qquad f_2=(0,b_2)
 \quad\hbox{in }H^0(V)
\]
gives multiplication--Cartier rank $6q-2$ and stable rank $6q-2$
on $E\times E_{\rm s}$, where
$E_{\rm s}:w^2+w=z^3$ and the second line bundle is $\OO(3O)$.
The corresponding section is $s=f_0+zf_1+wf_2$.
\end{lemma}
\begin{proof}
First choose $a_1,b_1$ in the nonempty open where
$(a_0,a_1)$ and $(b_1,b_2)$ are basepoint-free pencils.
These conditions say respectively that $a_1$ is nonzero on $A$
and $b_1$ is nonzero on $D$.
Write an arbitrary input triple in $H^0(V)^3$ as
$(x_i,y_i)$, with $x_i\in H^0(H)$ and $y_i\in H^0(HB)$.
On the supersingular factor the multiplication--Cartier pairing
in the basis $(1,z,w)$ is
\[
 \begin{pmatrix}0&1&z\\1&0&w\\z&w&0\end{pmatrix}.
\]
Consequently the linearized Cartier map $T_s$ has output
$g=(g_0,g_1,g_2)\in H^0(H^2)^3$ precisely when
\begin{align}
 a_0x_1+a_1x_0+\mu b_1y_0&=g_0,\label{eq:DE-residue-C0}\\
 a_0x_2+\mu b_2y_0&=g_1,\label{eq:DE-residue-C1}\\
 a_1x_2+\mu b_1y_2+\mu b_2y_1&=g_2.\label{eq:DE-residue-C2}
\end{align}
The middle equation determines $x_2,y_0$ uniquely, by the preceding
lemma.  The other two equations are the two pencil multiplication
maps, each with kernel and cokernel of dimension one.
In particular $T_s$ has kernel exactly
\[
 k\,k_a\oplus k\,k_b,\qquad
 k_a=((a_0,0),(a_1,0),(0,0)),\quad
 k_b=((0,0),(0,b_1),(0,b_2)),
\]
and its rank is $6q-2$.

Here is the precise Frobenius convention needed for stable rank.
Descent of squares gives a bijective Frobenius-semilinear map
\[
 Q:H^0(E,F^*H)\longrightarrow H^0(E',H^2).
\]
The identification $(F^*H)^{(2)}=H^2$ is canonical in the present
origin-divisor notation.  Componentwise,
$Q(a,b)=a^2+\mu b^2$.
The absolute semilinear Cartier operator is $Q^{-1}T_s$,
up to the harmless normalization of the invariant differential.
Thus its kernel meets its image trivially if and only if the map
\[
 Q(\ker T_s)\longrightarrow\coker T_s
\]
is an isomorphism.  The two output vectors to test are
\[
 Q(k_a)=(a_0^2,a_1^2,0),\qquad
 Q(k_b)=(0,\mu b_1^2,\mu b_2^2).
\]
This square-descent step prevents confusing a linear Cartier map
between Frobenius-twisted spaces with an ordinary endomorphism.

Fix a nonzero invariant differential $\omega$ on $E'$.
The two pencil cokernels, identified with $H^1(E',\OO)=k$, have
residue functionals
\[
 L_A(g)=\sum_{P\in A}\operatorname{Res}_P
              \frac{g\,\omega}{a_0a_1},\qquad
 L_D(g)=\sum_{P\in D}\operatorname{Res}_P
              \frac{g\,\omega}{\mu b_1b_2}.
\]
Indeed the two local solutions of each pencil equation differ by
the displayed fraction times its kernel generator.  This is the
two-open Cech connecting homomorphism followed by residue duality.
The original references are \cite[Theorem~2 and Theorem~3 with its
corollary, pp.~155--156; Section~4, Theorem~5 and its corollary,
pp.~157--159]{Tate1968Residues}.

In \eqref{eq:DE-residue-C1}, reduction at $A$ gives
$y_0=g_1/(\mu b_2)$, and reduction at $D$ gives $x_2=g_1/a_0$.
Subtracting these local substitutions changes the other expressions
only by functions regular at the respective residue points.
The two-obstruction matrix on $Q(k_a),Q(k_b)$ is therefore
\[
 \begin{pmatrix}t&u\\v&t\end{pmatrix},
\quad
\begin{split}
 t&=\sum_A\operatorname{Res}\frac{a_1b_1}{a_0b_2}\omega,\\
 u&=\sum_A\operatorname{Res}\frac{\mu b_1^3}{a_0a_1b_2}\omega,\\
 v&=\sum_D\operatorname{Res}\frac{a_1^3}{\mu a_0b_1b_2}\omega .
\end{split}
\]
For its lower-right entry, the same expression as $t$ is summed
over $D$.  The rational differential has poles only in $A\cup D$;
the global residue theorem and characteristic two make these sums
equal.  All displayed quotients are globally defined rational
functions: the line-bundle factors cancel using $\mu$.

Replace $b_1$ by $b_1+\lambda b_2$, with $\lambda\in k$.
The pencils remain basepoint-free, since the values on $D$ do not
change.  The entry $v$ is constant: at every simple zero of $b_2$
only the value of $b_1$ enters its residue.
The entry $t$ is constant as well, since its change is
\[
 \lambda\sum_A\operatorname{Res}\frac{a_1}{a_0}\omega=0.
\]
The entry $u$ is a polynomial of degree at most three in $\lambda$.
Its linear coefficient is
\[
 c_1=\sum_A\operatorname{Res}\frac{\mu b_1^2}{a_0a_1}\omega.
\]
For every fixed allowed $a_1$, this is a nonzero quadratic polynomial
in $b_1$.  To verify this without a genericity assumption on $E$,
use the evaluation isomorphism $H^0(HB)\simeq H^0(A,HB|_A)$.
At each reduced point of $A$ its coefficient is nonzero:
$a_0$ has a simple zero, and $a_1,\mu,\omega$ are units there.
The independent evaluation coordinates therefore exhibit $c_1$
as a nonzero sum of their squares with nonzero coefficients.

Similarly, for every fixed allowed $b_1$, the entry $v$ is a
nonzero cubic polynomial in $a_1$.  The evaluation map
$H^0(H)\simeq H^0(D,H|_D)$ is an isomorphism, and each of the
simple-pole residue coefficients on $D$ is nonzero.
Thus $v$ is a nonzero sum of cubes of independent evaluation
coordinates.  Both nonvanishing conditions define nonempty opens
in the irreducible open parameter space of allowed pairs $(a_1,b_1)$;
they have a common point.
Choose such a point.  Then $v\ne0$ and $u(\lambda)$ is nonconstant,
whereas $t$ and $v$ are constant.  Hence
\[
                  t^2+u(\lambda)v
\]
is a nonzero polynomial, and is nonzero for some $\lambda$.
The obstruction matrix is invertible for that choice.
Therefore the Cartier kernel has zero intersection with the Cartier
image.  Rank already stabilizes after the first iterate, and equals
$6q-2$, as claimed.
\end{proof}

\begin{theorem}[All mixed thin product systems]
\label{thm:DE-A3-mixed-thin-all}
\status{Proved; root and reduction-descent full independent proof/source/code audits passed.}
Let $E_{\rm o}$ be any ordinary elliptic curve and $E_{\rm s}$
any supersingular elliptic curve over an algebraically closed
field of characteristic two.  For every even $m\ge6$, a general
smooth member of
\[
 \bigl|\OO_{E_{\rm o}}(mO)\boxtimes\OO_{E_{\rm s}}(3O)\bigr|
\]
has genus $3m+1$, full stable rank $3m-1$, and relative stable
rank $3m-2$.
\end{theorem}
\begin{proof}
Write $m=2q$.  Over the algebraic closure the supersingular factor
has the displayed model used in the preceding lemma.
That lemma supplies a section for which the relative Cartier
stable rank is $3m-2$; it does not require that section itself
define a smooth curve.
The corresponding iterated-rank minor is nonzero on a nonempty
open of the complete section space.  Both line bundles have
degree at least three, so their exterior product is very ample.
Bertini supplies a nonempty smooth-member open, which intersects
the rank open.
Lemma~\ref{lem:DE-A3-mixed-even-bound} gives the matching generic
upper bound.  The divisor sequence, its stable-part exactness,
and Serre duality are precisely those proved there; the ambient
stable rank is one.  Adjunction gives genus $3m+1$.
The result is uniform in the ordinary elliptic curve and in $m$.
\end{proof}

\begin{corollary}[The mixed product formula in every degree]
\label{cor:DE-A3-mixed-product-no-exception}
\status{Proved; root and reduction-descent full independent proof/source/code audits passed.}
For every perfect field of characteristic two, every ordinary
$E_{\rm o}$ and supersingular $E_{\rm s}$, and all positive
integers $a,b,d$, the geometric generic relative stable rank in
\[
 \bigl|\OO_{E_{\rm o}}(3adO)\boxtimes
             \OO_{E_{\rm s}}(3bdO)\bigr|
\]
is
\[
            r_d=9abd^2-1-\epsilon(ad),\qquad
 \epsilon(c)=\begin{cases}1&c\ {\rm even},\\0&c\ {\rm odd}.
                         \end{cases}
\]
The relative coherent dimension is $9abd^2-1$.
\end{corollary}
\begin{proof}
Theorem~\ref{thm:DE-A3-mixed-all-product-ranks} already proves
every case except $d=b=1$ with $a$ even.
Apply Theorem~\ref{thm:DE-A3-mixed-thin-all} with $m=3a$
to the remaining cases.  Passage to the algebraic closure
computes the asserted geometric generic ranks over a perfect field.
The earlier ambient restriction bound remains applicable.
\end{proof}

\begin{remark}[Verification and scope of the residue argument]
\label{rem:DE-A3-mixed-thin-residue-checks}
\status{Proved; root and reduction-descent full independent proof/source/code audits passed.}
The entire new argument is uniform in $q$ and in the ordinary elliptic
curve.  Its two nonvanishing assertions use evaluation isomorphisms
on reduced divisors, rather than a list of elliptic parameters.
The exact symbolic program
\path{agents/obstructions_crystals/code/mixed_thin_residue_checks.py}
checks thirteen universal characteristic-two identities: the two
kernel vectors, the middle exact-sequence kernel, both regular
residue-remainder substitutions, all three terms of the cubic
variation, and the nonzero determinant coefficient $vc_1$.
The author run passed in $0.1971881189965643$ seconds, with its full
output retained beside the program.
The root and reduction-descent reviewers read the entire proof,
original source passages, and program; their author-preserving
reruns passed in $0.1764095999969868$ and
$0.19381124299980002$ seconds, respectively.
The ordinary Frobenius-algebra splitting and all global existence
steps are proved above, not inferred from that symbolic check.
This closes the thin-system exception for mixed elliptic products;
it does not supply a generic-rank formula for arbitrary ambient
varieties.
\end{remark}
